\SourceLeaf{248}%
With this choice, if we denote the variation by the symbol \(\delta\), we have
\[
\begin{aligned}
\delta s&=\int\frac{(\tfrac{\mathrm{d}E}{\mathrm{d}p}\mathord{.}\mathrm{d}p^2+\tfrac{2\mathrm{d}F}{\mathrm{d}p}\mathord{.}\mathrm{d}p\mathord{.}\mathrm{d}q+\tfrac{\mathrm{d}G}{\mathrm{d}p}\mathrm{d}q^2)\delta p+(2E\mathrm{d}p+2F\mathrm{d}q)\mathrm{d}\delta p}{2\mathrm{d}s}\\
&=\frac{E\mathrm{d}p+F\mathrm{d}q}{\mathrm{d}s}\mathord{.}\delta p+\int\delta p\mathord{.}\{\frac{\tfrac{\mathrm{d}E}{\mathrm{d}p}\mathord{.}\mathrm{d}p^2+\tfrac{2\mathrm{d}F}{\mathrm{d}p}\mathord{.}\mathrm{d}p\mathord{.}\mathrm{d}q+\tfrac{\mathrm{d}G}{\mathrm{d}p}\mathord{.}\mathrm{d}q^2}{2\mathrm{d}s}-\mathrm{d}\mathord{.}\frac{E\mathrm{d}p+F\mathrm{d}q}{\mathrm{d}s}\}
\end{aligned}
\]
It is clear that what stands here under the integral sign must vanish independently of \(\delta p\). We therefore obtain
\[
\begin{aligned}
\frac{\mathrm{d}E}{\mathrm{d}p}\mathord{.}\mathrm{d}p^2+\frac{2\mathrm{d}F}{\mathrm{d}p}\mathord{.}\mathrm{d}p\mathord{.}\mathrm{d}q+\frac{\mathrm{d}G}{\mathrm{d}p}\mathord{.}\mathrm{d}q^2&=2\mathrm{d}s\mathord{.}\mathrm{d}\mathord{.}\frac{E\mathrm{d}p+F\mathrm{d}q}{\mathrm{d}s}\\
&=2\mathrm{d}s\mathord{.}\mathrm{d}\mathord{.}\BareRoot{E}\mathord{.}\cos\theta\\
&=\frac{\mathrm{d}s\mathord{.}\mathrm{d}E\mathord{.}\cos\theta}{\BareRoot{E}}-2\mathrm{d}s\mathord{.}\mathrm{d}\theta\mathord{.}\BareRoot{E}\mathord{.}\sin\theta\\
&=\frac{(E\mathrm{d}p+F\mathrm{d}q)\mathrm{d}E}{E}-\BareRoot{(EG-FF)}\mathord{.}\mathrm{d}q\mathord{.}\mathrm{d}\theta\\
&=(\frac{E\mathrm{d}p+F\mathrm{d}q}{E})\mathord{.}(\frac{\mathrm{d}E}{\mathrm{d}p}\mathord{.}\mathrm{d}p+\frac{\mathrm{d}E}{\mathrm{d}q}\mathord{.}\mathrm{d}q)\\
&\quad-2\BareRoot{(EG-FF)}\mathord{.}\mathrm{d}q\mathord{.}\mathrm{d}\theta
\end{aligned}
\]
This gives us the following equation expressing the condition for a shortest line:
\[
\begin{aligned}
\BareRoot{(EG-FF)}\mathord{.}\mathrm{d}\theta
&=\tfrac{1}{2}\frac{F}{E}\mathord{.}\frac{\mathrm{d}E}{\mathrm{d}p}\mathord{.}\mathrm{d}p+\tfrac{1}{2}\mathord{.}\frac{F}{E}\mathord{.}\frac{\mathrm{d}E}{\mathrm{d}q}\mathord{.}\mathrm{d}q\\
&\quad+\tfrac{1}{2}\mathord{.}\frac{\mathrm{d}E}{\mathrm{d}q}\mathord{.}\mathrm{d}p-\frac{\mathrm{d}F}{\mathrm{d}p}\mathrm{d}p-\tfrac{1}{2}\mathord{.}\frac{\mathrm{d}G}{\mathrm{d}p}\mathrm{d}q
\end{aligned}
\]
which can also be written
\[
\BareRoot{(EG-FF)}\mathord{.}\mathrm{d}\theta=\tfrac{1}{2}\frac{F}{E}\mathord{.}\mathrm{d}E+\tfrac{1}{2}\mathord{.}\frac{\mathrm{d}E}{\mathrm{d}q}\mathord{.}\mathrm{d}p-\frac{\mathrm{d}F}{\mathrm{d}p}\mathord{.}\mathrm{d}p-\tfrac{1}{2}\mathord{.}\frac{\mathrm{d}G}{\mathrm{d}p}\mathord{.}\mathrm{d}q
\]
Furthermore, by means of the equation
\[
\mathrm{cotg}\,\theta=\frac{E}{\BareRoot{(EG-FF)}}\mathord{.}\frac{\mathrm{d}p}{\mathrm{d}q}+\frac{F}{\BareRoot{(EG-FF)}}
\]
we can eliminate the angle \(\theta\) from that equation and thereby derive a second-order differential equation relating \(p\) and \(q\). It would, however, be more complicated and less useful for applications than the preceding one.
\Article{19}
The general formulas that we derived in Articles 11 and 18 for the measure of curvature and the change in direction of a shortest line become much simpler if the quantities \(p,q\) are chosen so that the lines of the first system intersect those of the second at right angles everywhere, that is, so that \(\omega=90^0\) or \(F=0\) holds in general. The formula for the measure of curvature then becomes
